% Solutions of Book 4, chapter 16 (sockets).
\section*{Chapter \ref{ch:b2:sockets} --- Sockets}

\begin{solution}{exo:b2:sockets:1}
The client receives the bytes \code{abcd} in order, cut anywhere into pieces of at most 3 bytes, each a nonempty
prefix of what is left. The possible outputs are the compositions of 4 into parts of at most 3: \code{[abc][d]},
\code{[a][bcd]}, \code{[ab][cd]}, \code{[ab][c][d]}, \code{[a][bc][d]}, \code{[a][b][cd]} and \code{[a][b][c][d]},
7 outputs. The two sends do not force a cut between \code{b} and \code{c}: \code{[abc][d]} happens if both had
arrived before the first \code{recv}.
\end{solution}

\begin{solution}{exo:b2:sockets:2}
$44164 = \mathtt{0xac84}$. Stored without \code{htons} on this little-endian machine, \code{sin_port} holds the bytes
\code{84 ac}; the system reads them as a big-endian port, $\mathtt{0x84ac} = 33\,964$, as \code{getsockname}
reports. \code{bind} succeeded because 33\,964 is an ordinary free port: nothing checks that it is the one meant. On a
big-endian machine the bug would be invisible, as \code{htons} does nothing there; that is why the conversion must be
written even where it is not needed.
\end{solution}

\begin{solution}{exo:b2:sockets:3}
\omcode{code/bachelor-2/16-sockets/listings/net.c}{44}{54}{Sending everything}
\code{recv_all} is \cref{lst:b2:sockets:recv-all}, proved by \cref{prop:b2:sockets:recv-all}. \code{send_all}
loops the same way: \code{send} may accept fewer bytes than given (when the \omterm{def:b2:sockets:socket}{socket}'s send buffer is nearly full), so
the rest is sent again from where it stopped. \code{MSG_NOSIGNAL} turns the \code{SIGPIPE} of a write to a closed
connection into an error return. The test sends a frame and lines one byte at a time and checks that they arrive
whole.
\end{solution}

\begin{solution}{exo:b2:sockets:4}
\code{street Mill Lane} has 16 bytes, so the length is \code{00 00 00 10}, followed by \code{73 74 72 65 65 74 20 4d
69 6c 6c 20 4c 61 6e 65}: 20 bytes in all. \code{recv_msg} reads them in two calls to \code{recv_all}, one for the 4
bytes of the length and one for the 16 of the message; each of those may make several \code{recv} calls.
\end{solution}

\begin{solution}{exo:b2:sockets:5}
\omoutput{code/bachelor-2/16-sockets/outputs/exo_truncate.txt}{The first datagram is cut; the second is whole.}
A \code{recvfrom} returns one datagram. When it is larger than the buffer, the first 10 bytes are copied and the other
90 are discarded with the datagram: they are not kept for the next call, which returns the next datagram,
\code{next}. A program that cares passes \code{MSG_TRUNC}, with which Linux returns the datagram's real length, or
uses a buffer as large as the largest datagram its protocol allows.
\end{solution}

\begin{solution}{exo:b2:sockets:6}
After \code{hello} and its echo, the client closes. Every later \code{recv} returns 0 at once, end of stream, and the loop goes
round without waiting: 999\,999 calls in a fraction of a second, all returning 0, burning a processor. Without the
guard it would never stop. The fix is to leave the loop on 0 and on $-1$: \code{while ((n = recv(c, buf, sizeof buf,
0)) > 0) send_all(c, buf, n);}, then \code{close(c)}. Sending back \code{n} bytes when \code{n} is $-1$ would be a
second bug, which the guard's \code{n > 0 ? n : 0} hides.
\end{solution}

\begin{solution}{exo:b2:sockets:7}
\omcode{code/bachelor-2/16-sockets/exercises/linereader.c}{13}{37}{A line reader over a stream}
The unread bytes live in \code{buf[start..end)}. If they contain a newline, the line is returned and \code{start}
moves past it, keeping the bytes after it for the next call: one \code{recv} may bring several lines. Otherwise the
rest is moved to the front and \code{recv} fills the free space: a line may come in pieces. A line longer than the
buffer is dropped, and one longer than \code{cap} is cut; a final line without newline is lost at the end of the
stream, a choice to document. The test sends lines one byte at a time, including an empty line.
\end{solution}

\begin{solution}{exo:b2:sockets:8}
By induction on $k$. For $k = 0$, $w$ is empty, and decoding yields nothing and uses all of it. Let $k \ge 1$ and $w =
f(m_1) w'$, where $f(m_1)$ is 4 bytes holding $n_1 = |m_1|$ followed by $m_1$, and $w'$ is the concatenation of the
framings of $m_2, \dots, m_k$. Decoding reads the first 4 bytes of $w$, which are those of $f(m_1)$, so it reads
$n_1$; it then takes the next $n_1$ bytes, which are $m_1$, and starts again at the first byte of $w'$. By the
induction hypothesis, decoding $w'$ yields $m_2, \dots, m_k$ and uses all of $w'$. So decoding $w$ yields $m_1, \dots,
m_k$ and uses all of it. The key point is that the decoder never has to guess where a message ends: the length says
so before any byte of the message is read.

With the newline as delimiter, the single message made of \code{a}, a newline and \code{b}, and the two messages
\code{a} and \code{b}, are both sent as the four bytes \code{a}, newline, \code{b}, newline: the receiver cannot tell
them apart.
\end{solution}

\begin{solution}{exo:b2:sockets:9}
\omcode{code/bachelor-2/16-sockets/exercises/count_server.c}{16}{32}{One \omterm{def:b2:threads-and-synchronisation:thread}{thread} per client, one shared
counter}
The count is shared by every \omterm{def:b2:threads-and-synchronisation:thread}{thread}, so its increment and the read of its new value are one \omterm{def:b2:threads-and-synchronisation:mutex-property}{critical section}; reading
\code{lines} again after unlocking could return another \omterm{def:b2:threads-and-synchronisation:thread}{thread}'s number. The reply is formatted and sent outside the
\omterm{def:b2:threads-and-synchronisation:lock}{lock}, so that a slow client does not hold the others up. The line reader of \cref{exo:b2:sockets:7} cuts the stream
into lines. The test runs 4 clients of 50 lines each and checks that the server counted 200, and that no reply
exceeds 200.
\end{solution}

\begin{solution}{exo:b2:sockets:10}
\omcode{code/bachelor-2/16-sockets/exercises/chat.c}{21}{57}{A chat server: one \code{poll}, one buffer per client}
Each client has a buffer in which partial lines wait; every complete line is sent to every other client, and removed.
A client that closes, or fails, is closed and forgotten, its slot taken by the last one (the loop runs downwards so
that the moved entry has already been handled). The \omterm{def:b2:sockets:listening}{listening socket} is in the set only while clients may still
come. \code{send_all} may block if a client stops reading; a server for real use would keep an outgoing buffer per
client and ask \code{poll} for \code{POLLOUT}. The test connects three clients, sends a line from one and checks
that the two others receive it, then the same from another.
\end{solution}

\begin{solution}{exo:b2:sockets:11}
Eight clients keep the server busy, so the time per connection is the inverse of the rate. On the run printed, the
\omterm{def:b2:sockets:event-loop}{event loop} takes $1/10\,917 \approx 91.6$~µs per connection, which we take as $r$. The \omterm{def:b2:threads-and-synchronisation:thread}{thread} design takes
$1/7\,331 \approx 136.4$~µs, so $f \approx 45$~µs for a \omterm{def:b2:threads-and-synchronisation:thread}{thread}; the \omterm{def:b2:processes-and-system-calls:process}{process} design $1/3\,662 \approx 273.1$~µs, so
$f \approx 180$~µs for a \omterm{def:b2:processes-and-system-calls:process}{process}. With $q$ requests per connection, a design serves $q / (f + q r)$ requests per unit
of time, against $1/r$ for the \omterm{def:b2:sockets:event-loop}{event loop}. The \omterm{def:b2:processes-and-system-calls:process}{process} design reaches 90~\% of it when $q r \ge 0.9 (f + q r)$, that
is $q \ge 9 f / r \approx 9 \times 180 / 91.6 \approx 18$ requests per connection. The model ignores everything else,
such as the eight client \omterm{def:b2:threads-and-synchronisation:thread}{threads} competing for the same cores; it says that the fixed cost of a design matters only
for short connections.
\end{solution}

\begin{solution}{exo:b2:sockets:12}
The side that closes a TCP connection first keeps it, after the exchange of closing messages, in the state
\code{TIME_WAIT} for a while, so that a late packet of the old connection cannot be taken for part of a new one with
the same addresses and ports. Here the server closed first, so its end, bound to port 44\,169, is in
\code{TIME_WAIT} after the program has closed every descriptor. A new \code{bind} to that port then fails with
\code{EADDRINUSE}, unless \code{SO_REUSEADDR} is set, which allows binding an address whose only users are such
lingering connections (it still refuses while another \omterm{def:b2:sockets:socket}{socket} is listening on it). This is why a server restarted at
once fails without the option, and why every server of the chapter sets it.
\end{solution}

\begin{solution}{pb:b2:sockets:1}
\textbf{1.} \code{00 00 00 05 77 68 6f 20 31}: the length 5, then \code{w}, \code{h}, \code{o}, a space, \code{1}.

\textbf{2.} $2^{32} - 1$ bytes, about 4 GiB. A malicious or broken client could announce that much; the capacity
bounds what the server agrees to receive, and \code{recv_msg} returns $-1$ for a longer message instead of
overflowing its buffer.

\textbf{3.} Yes: the stream keeps no boundaries, and the first \code{recv} may return all 18 bytes. Framing makes it
harmless, because the decoder takes exactly the announced length and leaves the rest for the next message
(\cref{prop:b2:sockets:prefix-free}).

\textbf{4.} 39 = \code{0x27}: \code{00 00 00 27}.

\textbf{5.} \code{"<I"} is little-endian: \code{05 00 00 00}. The server would read $\mathtt{0x05000000} = 83\,886\,080$
as the length, more than its capacity, and \code{recv_msg} would return $-1$: the server would close without an
answer.

\textbf{6.} A closed connection says only that something went wrong; an error message says what, and lets the client
tell a mistake in its request from a broken server or network.

\textbf{7.} \code{socket} creates a \omterm{def:b2:sockets:kinds}{stream socket}; \code{setsockopt} sets \code{SO_REUSEADDR}; \code{bind} gives it
the address 127.0.0.1 and the port; \code{listen} makes it a \omterm{def:b2:sockets:listening}{listening socket} with a queue of 128 pending connections.

\textbf{8.} A new descriptor for the one connection with that client. The \omterm{def:b2:sockets:listening}{listening socket} goes on receiving other
clients' connections; each connection needs its own descriptor to send and receive on.

\textbf{9.} \code{recv_msg} waits for 4 bytes. With a \omterm{def:b2:processes-and-system-calls:process}{process} or a \omterm{def:b2:threads-and-synchronisation:thread}{thread} per connection, only that child or \omterm{def:b2:threads-and-synchronisation:thread}{thread}
waits, the others go on. In the \omterm{def:b2:sockets:event-loop}{event loop}, \code{poll} never reports the silent connection ready, so the loop goes on
too; the connection simply stays in the set (until the client leaves).

\textbf{10.} \code{recv_all(fd, len, 4)} returns 4 with the length, or fewer at the end of the stream, and
\code{recv_msg} returns $-1$ unless it got 4. Then \code{recv_all(fd, buf, n)} returns $n$ with the message, or fewer,
and again \code{recv_msg} returns $-1$ unless it got $n$. So it returns $n$ with a whole message, or $-1$.

\textbf{11.} A write to a connection the client has closed raises \code{SIGPIPE}, whose default action ends the
\omterm{def:b2:processes-and-system-calls:process}{process}: the whole server would die because of one client. With \code{MSG_NOSIGNAL}, \code{send} fails with
\code{EPIPE} and only that connection is closed.

\textbf{12.} It is told to: its third argument is the number of connections to serve, 4 here, one per request of the
client. A server that never stops would leave the session waiting.

\textbf{13.} The parent's copy keeps the connection open: the client would see end of stream only when both copies
are closed, and each such connection would hold a descriptor in the parent until it ran out of them.

\textbf{14.} \code{&c} points to a variable that the next turn of the loop overwrites, perhaps before the \omterm{def:b2:threads-and-synchronisation:thread}{thread} has
read it: two \omterm{def:b2:threads-and-synchronisation:thread}{threads} could serve the same descriptor. The value itself, carried in the pointer, belongs to the
\omterm{def:b2:threads-and-synchronisation:thread}{thread}.

\textbf{15.} The tables are written once, before the server starts, and only read afterwards. Readers alone do not
form a \omterm{def:b2:threads-and-synchronisation:data-race}{data race}.

\textbf{16.} The first piece is appended to \code{st[i].buf}; \code{unframe} finds the message incomplete and the
connection stays in the set. When the second piece arrives, \code{poll} reports it again, and the request is then
complete.

\textbf{17.} Removing entry $i$ moves the last entry into slot $i$. Walking downwards, that last entry has already been
handled in this round; walking upwards, it would be skipped in this round (the copy carries its own \code{revents}, but the loop has
passed slot $i$), and since \code{poll} reports a ready descriptor again, it would be served only after the next
call. Walking downwards serves every ready entry at once.

\textbf{18.} The \omterm{def:b2:processes-and-system-calls:process}{process} per connection: a crash kills only that child. A \omterm{def:b2:threads-and-synchronisation:thread}{thread}'s crash kills the whole \omterm{def:b2:processes-and-system-calls:process}{process}, every
\omterm{def:b2:threads-and-synchronisation:thread}{thread} with it; the \omterm{def:b2:sockets:event-loop}{event loop} is one \omterm{def:b2:threads-and-synchronisation:thread}{thread}, so it dies too.

\textbf{19.} $8 \times 500 = 4000$ connections, from 8 client \omterm{def:b2:threads-and-synchronisation:thread}{threads}.

\textbf{20.} The server must be measured as a separate program, as it would run: in a \omterm{def:b2:processes-and-system-calls:tree}{child process}, its \code{fork}
copies only the small server, not the client \omterm{def:b2:threads-and-synchronisation:thread}{threads}, and a crash or a stuck \omterm{def:b2:threads-and-synchronisation:thread}{thread} of one side does not take the other
with it.

\textbf{21.} The clients see the whole service time, from the first connection to the last answer, which is what a
user experiences; the server cannot know when the clients are done.

\textbf{22.} \omterm{def:b2:sockets:event-loop}{Event loop}, \omterm{def:b2:threads-and-synchronisation:thread}{thread} per connection, \omterm{def:b2:processes-and-system-calls:process}{process} per connection; on the run printed the \omterm{def:b2:sockets:event-loop}{event loop} is about
$10\,917 / 3\,662 \approx 3.0$ times as fast as the \omterm{def:b2:processes-and-system-calls:process}{process} design.

\textbf{23.} The fixed cost per connection would be paid once for a thousand requests, so the three designs would come
close (\cref{exo:b2:sockets:11}); the \omterm{def:b2:threads-and-synchronisation:thread}{threads} and \omterm{def:b2:processes-and-system-calls:process}{processes} might even pull ahead, using several cores where the \omterm{def:b2:sockets:event-loop}{event
loop} uses one.

\textbf{24.} Requests per second on loopback for three server designs, on the laptop: about 3\,700 with a \omterm{def:b2:processes-and-system-calls:process}{process} per
connection, 7\,300 with a \omterm{def:b2:threads-and-synchronisation:thread}{thread} per connection and 10\,900 with an \omterm{def:b2:sockets:event-loop}{event loop}, in the run printed
(3662, 7331 and 10\,917).
\end{solution}
